\chapter{Le carnet d'ordres à cours limité}\label{ch:mx:the-limit-order-book}

Un trader envoie un ordre d'achat de 100 actions au meilleur cours acheteur. L'écran affiche un prix et une taille~; ce qu'il n'affiche pas, c'est une file. Dans
l'heure simulée que ce livre utilise d'un bout à l'autre, 2\,834 ordres se sont joints à l'arrière du meilleur cours acheteur ou vendeur~: ils ont trouvé devant
eux une médiane de 1\,400 actions, et seuls 13~\% d'entre eux ont fini par être exécutés. Le meilleur prix n'est pas un lieu où un ordre repose~; c'est une file dans
laquelle un ordre attend, et l'ordre dans lequel la file s'est formée détermine qui échange. Ce chapitre décrit le carnet comme un état, comme un flux d'événements qui modifient cet état, comme une
structure de données qu'une plateforme et chaque participant doivent tenir à jour, et comme un système de files d'attente dont le modèle le plus simple indique déjà combien de temps dure l'attente.

\section{État~: niveaux de prix et files}

\begin{definition}[Ordre à cours limité, ordre au marché]\label{def:mx:the-limit-order-book:orders}
Un \emph{ordre à cours limité}\index{ordre à cours limité} est un ordre d'achat à un prix donné ou inférieur, ou de vente à un prix donné ou supérieur~; la part
qui ne peut pas s'exécuter immédiatement reste dans le carnet jusqu'à son exécution, son annulation ou son expiration. Un \emph{ordre au marché}\index{ordre au marché} est un ordre
d'achat ou de vente d'une quantité aux meilleurs prix disponibles~; il ne reste jamais au carnet.
\end{definition}

\begin{definition}[Pas de cotation, niveau de prix]\label{def:mx:the-limit-order-book:tick}
Le \emph{pas de cotation}\index{pas de cotation} $\delta_{\mathrm{tick}}$ est la plus petite variation de prix qu'une plateforme accepte pour un instrument~: chaque prix limite
est un multiple entier de celui-ci. Un \emph{niveau de prix}\index{niveau de prix} est l'ensemble des ordres au repos d'un même côté du carnet, à un même prix.
\end{definition}

\begin{definition}[Carnet d'ordres à cours limité]\label{def:mx:the-limit-order-book:lob}
Un \emph{carnet d'ordres à cours limité}\index{carnet d'ordres à cours limité} est l'ensemble des \omterm{def:mx:the-limit-order-book:orders}{ordres à cours limité} au repos portant sur un instrument donné sur une plateforme donnée, organisés par côté et par
\omterm{def:mx:the-limit-order-book:tick}{niveau de prix}~; à l'intérieur d'un niveau, par priorité. Son meilleur cours acheteur $b_t$ et son meilleur cours vendeur $a_t$ sont le plus haut prix d'achat et le plus bas prix de vente au repos~;
$Q^b_t$ et $Q^a_t$ sont les quantités affichées à ces prix.
\end{definition}

Un carnet est \emph{croisé} si $b_t\ge a_t$~; un carnet en négociation continue ne l'est jamais, car un ordre qui croiserait s'exécute à la place. La
fourchette $s_t = a_t - b_t$ vaut au moins un tick. Trois vues du même état reviennent dans tous les chapitres suivants (livre~1, chapitre~28)~: le niveau 1 est
$(b_t, Q^b_t, a_t, Q^a_t)$~; le niveau 2 ajoute la taille agrégée à chaque prix plus profond~; le niveau 3 est le carnet lui-même, chaque ordre avec sa référence, dans
l'ordre de la file (\cref{fig:mx:the-limit-order-book:book}).

\begin{omfigure}
\begin{tikzpicture}[font=\small,x=1cm,y=0.62cm]
\draw[->,thick] (0,0) -- (0,7.4) node[above]{prix};
\foreach \p/\lab in {1/99.97,2/99.98,3/99.99,4/100.00,5/100.01,6/100.02,7/100.03} {\draw (-0.08,\p) -- (0.08,\p); \node[left,font=\footnotesize] at (-0.1,\p) {\lab};}
\foreach \y/\x/\w in {4/0.4/1.2,4/1.65/0.8,4/2.5/1.6,3/0.4/2.0,3/2.45/1.0,2/0.4/1.4,2/1.85/0.6,2/2.5/1.1,1/0.4/2.4} {\draw[fill=omDef!18,draw=omDef] (\x,\y-0.32) rectangle ++(\w,0.64);}
\foreach \y/\x/\w in {5/0.4/0.9,5/1.35/1.5,6/0.4/1.8,6/2.25/0.7,7/0.4/1.3,7/1.75/1.3} {\draw[fill=omThm!18,draw=omThm] (\x,\y-0.32) rectangle ++(\w,0.64);}
\draw[fill=omExo!45,draw=omExo,thick] (4.15,3.68) rectangle ++(0.7,0.64); \node[font=\footnotesize] at (4.5,4) {achat};
\node[anchor=west,font=\footnotesize,align=left] at (5.0,4) {nouvel achat~: s'ajoute à l'arrière\\du meilleur cours acheteur};
\draw[->,thick,omThm] (5.8,5.52) -- (0.85,5.36) node[pos=0,right,font=\footnotesize,align=left]{un achat au marché prend\\le meilleur cours vendeur\\par l'avant};
\draw[decorate,decoration={brace,amplitude=4pt}] (-1.25,4.05) -- (-1.25,4.95) node[midway,left=5pt,font=\footnotesize]{fourchette};
\node[font=\footnotesize,anchor=west,omDef] at (0.4,-0.6) {achats~: priorité de prix vers le bas, puis temps};
\node[font=\footnotesize,anchor=west,omThm] at (1.2,7.75) {ventes~: priorité de prix vers le haut, puis temps};
\end{tikzpicture}

\omcaption{Un \omterm{def:mx:the-limit-order-book:lob}{carnet d'ordres à cours limité} au niveau 3. Chaque boîte est un ordre au repos, sa largeur étant sa taille~; à l'intérieur d'un \omterm{def:mx:the-limit-order-book:tick}{niveau de prix}, l'ordre le plus ancien
est à gauche. Le meilleur cours acheteur vaut 100.00 et le meilleur cours vendeur 100.01, à un tick d'écart. Un nouvel achat à 100.00 s'ajoute à l'arrière de son niveau~; un achat au marché consomme le
meilleur cours vendeur par l'avant.}\label{fig:mx:the-limit-order-book:book}
\end{omfigure}

\section{Événements~: ajout, annulation, modification, exécution}

\begin{definition}[Événement du carnet]\label{def:mx:the-limit-order-book:event}
Un \emph{événement du carnet}\index{evenement@événement du carnet} est une modification unique du carnet~: un ordre ajouté à un \omterm{def:mx:the-limit-order-book:tick}{niveau de prix}, réduit ou retiré par son propriétaire
(une annulation), remplacé par un nouvel ordre, ou réduit ou retiré par une exécution contre un ordre entrant. Le flux market-by-order d'une plateforme
(livre~1, chapitre~19) publie un message par événement, avec la référence de l'ordre.
\end{definition}

Le tableau associe les événements aux messages du flux TotalView-ITCH 5.0 de Nasdaq, dont le simulateur de bourse
du chapitre~26 reprend les formats. La spécification de Nasdaq est explicite sur le remplacement~: toutes les actions restantes de l'ordre initial sont retirées, et le
remplacement arrive avec une nouvelle référence d'ordre, utilisée dès lors, ce qui signifie pour un flux que l'ordre a perdu sa place.

\begin{center}
\small
\begin{tabular}{@{}llp{4.3cm}p{4.3cm}@{}}
\toprule
événement & message ITCH & niveau 3 & niveaux 2 et 1 \\
\midrule
ajout & A (add order) & un nouvel ordre à l'arrière de son niveau & la taille du niveau augmente~; le meilleur peut changer \\
annulation partielle & X (order cancel) & l'ordre diminue, garde sa place & la taille du niveau baisse \\
annulation totale & D (order delete) & l'ordre quitte le carnet & le niveau peut disparaître \\
remplacement & U (order replace) & l'ancien ordre quitte, le nouveau se place à l'arrière & deux niveaux peuvent changer \\
exécution & E / C (executed) & l'ordre de tête diminue ou quitte & le meilleur niveau diminue~; une transaction est publiée \\
\bottomrule
\end{tabular}
\end{center}

Les exécutions sont rares. Dans l'heure simulée (\texttt{firm.tape}, le marché synthétique du livre~7, 31\,535 messages), 49.7~\% des messages sont des ajouts,
47.8~\% des annulations et 2.5~\% des exécutions~: 40 messages pour chaque transaction. Seuls 3.8~\% des ordres qui ont quitté le carnet l'ont fait par une transaction~; la médiane
de durée de vie d'un ordre est de 15.5 secondes (24.9 pour ceux qui ont échangé, 15.3 pour ceux qui ont été annulés), et le meilleur cours acheteur a tenu une médiane de 4.5 secondes avant
de bouger. Les carnets réels sur les actions liquides vont plus vite et annulent davantage (le chapitre~3 place les faits simulés à côté de faits publiés)~; ce qui compte ici est la
proportion~: le carnet est surtout composé d'ordres qui n'échangeront jamais.

\section{Priorité et ce qui la remet à zéro}

Un \omterm{def:mx:the-limit-order-book:engine}{moteur d'appariement} (\cref{def:mx:the-limit-order-book:engine}) classe les ordres au repos par prix, puis selon une règle à l'intérieur du niveau~; la priorité prix-temps,
premier arrivé premier servi, est la règle courante pour les actions (livre~1, chapitre~19, présente les alternatives au prorata et à l'ordre de tête des
plateformes de contrats à terme). La priorité est une propriété de l'\emph{arrivée} d'un ordre, et les plateformes doivent préciser quelles modifications comptent comme une nouvelle arrivée. Les règles sont proches
d'une plateforme à l'autre~:
\begin{itemize}
\item \textbf{Une réduction conserve la place.} Une annulation partielle (ITCH X) laisse l'ordre où il se trouvait~: elle ne retire que ses propres actions.
\item \textbf{Une augmentation ou un changement de prix la fait perdre.} CME Globex énumère les modifications qui changent la priorité d'un ordre~: une augmentation de la quantité,
le prix, le type d'ordre, le prix de déclenchement du stop, la quantité minimale d'exécution et le compte. Sur Nasdaq, un cancel-replace arrive dans le flux sous
forme d'un Order Replace avec une nouvelle référence.
\item \textbf{Une tranche d'iceberg rafraîchie passe à l'arrière} (chapitre~2)~: montrer davantage d'une réserve cachée est une augmentation.
\end{itemize}
La raison en est l'équité envers les ordres placés derrière~: si une augmentation de taille conservait sa place, un trader pourrait se tenir en tête avec un seul lot et l'agrandir
dès qu'un ordre important arrive. Pour un trader, la conséquence est qu'une position dans la file est un actif (le chapitre~6 le valorise) qu'une modification mal conçue
détruit.

\begin{definition}[Moteur d'appariement]\label{def:mx:the-limit-order-book:engine}
Un \emph{moteur d'appariement}\index{moteur d'appariement} est le programme de la plateforme qui tient les \omterm{def:mx:the-limit-order-book:lob}{carnets d'ordres à cours limité}, traite chaque message entrant
dans l'ordre où il arrive, exécute les ordres entrants contre les ordres au repos selon la règle de priorité, et publie les événements et les rapports d'exécution qui en résultent.
\end{definition}

\section{Le carnet comme structure de données}

Un carnet comporte trois opérations sur son chemin critique~: retrouver un ordre par sa référence (chaque annulation et chaque exécution la nomme), trouver le meilleur prix (chaque
ordre agressif y commence), et parcourir un niveau dans l'ordre de priorité (chaque exécution). \texttt{firm.lob}, le carnet de référence du projet en cours,
conserve une table de hachage qui associe une référence à un ordre, une table qui associe un prix à un niveau de chaque côté, et deux files premier entré premier sorti par niveau, les ordres affichés
en premier (\cref{fig:mx:the-limit-order-book:ds}).

\begin{omfigure}
\begin{tikzpicture}[font=\footnotesize,box/.style={draw,rounded corners=2pt,minimum height=0.55cm,inner sep=3pt}]
\node[box,fill=omDef!12,minimum width=2.5cm,align=center] (h) at (0,2.6) {table de hachage\\ référence $\to$ ordre};
\node[box,fill=omDef!12,minimum width=2.5cm,align=center] (m) at (0,0) {cours acheteur~: prix $\to$ niveau\\ (trié)};
\foreach \i/\p in {0/100.00,1/99.99,2/99.98} {\node[box,minimum width=1.3cm] (p\i) at (3.3,0.9-0.9*\i) {\p};}
\foreach \k in {0,1,2} {\node[box,fill=omDef!18] (o\k) at (5.2+1.4*\k,0.9) {\#\the\numexpr 11+\k\relax};}
\draw[->] (p0) -- (o0); \draw[->] (o0) -- (o1); \draw[->] (o1) -- (o2);
\node[box,fill=gray!15] (hid) at (9.6,0.9) {caché};
\draw[->] (o2) -- (hid);
\draw[->] (m.east) -- (p0.west); \draw[->] (m.east) -- (p1.west); \draw[->] (m.east) -- (p2.west);
\draw[->,dashed,omDef] (h.east) to[out=0,in=100] (o1.north);
\node[anchor=west,align=left] at (5.0,-0.6) {niveau 100.00~: file affichée dans l'ordre d'arrivée,\\ puis la file cachée};
\end{tikzpicture}

\omcaption{Les structures de \texttt{firm.lob}. Une annulation retrouve son ordre grâce à la table de hachage (en pointillés)~; un ordre agressif trouve le meilleur
niveau grâce à la table triée et parcourt ses files depuis leur tête.}\label{fig:mx:the-limit-order-book:ds}
\end{omfigure}

\omcode{code/firm/lob/firm_lob.py}{60}{91}{Le cœur de \texttt{firm.lob}~: ajout, recherche, suppression et réduction.}\label{lst:mx:the-limit-order-book:lob}

\begin{proposition}[Coût des opérations sur le carnet]\label{prop:mx:the-limit-order-book:cost}
Avec $L$ niveaux non vides d'un côté et $m$ ordres à un niveau, la structure de la \cref{fig:mx:the-limit-order-book:ds} trouve un ordre en temps moyen
$O(1)$, ajoute un ordre en $O(\log L)$ lorsqu'il crée un niveau et en $O(1)$ sinon, trouve le meilleur prix en $O(1)$ (l'extrémité de la table triée), et
supprime un ordre en $O(\log L+m)$ dans le pire cas. Un tableau indexé par le prix couvrant $W$ ticks, muni d'un pointeur vers le meilleur indice, ajoute et supprime en $O(1)$
et déplace le pointeur du meilleur indice en $O(g)$, où $g$ est le nombre de ticks vides qu'il franchit.
\end{proposition}

\begin{proof}
La recherche dans la table de hachage est de coût constant en moyenne. Une table triée (un arbre équilibré, ou une liste triée avec recherche dichotomique, comme dans la référence Python) insère une nouvelle clé
en $O(\log L)$. La suppression localise l'ordre dans sa file, en temps linéaire en $m$ sauf si chaque ordre conserve un lien vers ses voisins (une liste doublement chaînée intrusive,
le choix de production, le rend en $O(1)$). Dans le tableau, un niveau est un indice calculé à partir du prix~; après qu'une suppression a vidé le meilleur niveau, le
pointeur passe à l'indice non vide suivant, soit $g$ pas.
\end{proof}

Les carnets de production pour les instruments liquides utilisent le tableau~: les prix se concentrent à quelques centaines de ticks du meilleur, $g$ est petit, et chaque
opération accède à la mémoire à une adresse calculée (le livre~13 en construit un). \texttt{firm.lob} réunit les deux~: son \texttt{ArrayBook} et son
\texttt{MessageBook}, à clés en arbre, a reconstruit le carnet de l'heure simulée message par message et concorde sur les 3\,154 instantanés de niveau 2 comparés.
Ses jumeaux en C++20 et en Rust reconstruisent un jeu de référence de cinq minutes et reproduisent chacun de ses 147 instantanés.

\section{Le carnet comme système de files d'attente}

Vu d'un seul ordre, le meilleur niveau est une file dans laquelle il attend~: les ordres qui le précèdent sortent par exécution (depuis la tête) ou par annulation (depuis
n'importe où), et l'ordre lui-même peut être annulé avant que son tour n'arrive. Le modèle le plus simple de l'ensemble du carnet traite chaque flux d'ordres comme aléatoire.

\begin{definition}[Modèle à intelligence nulle]\label{def:mx:the-limit-order-book:zi}
Un \emph{modèle à intelligence nulle}\index{modèle à intelligence nulle} (zero-intelligence model) d'un \omterm{def:mx:the-limit-order-book:lob}{carnet d'ordres à cours limité} fait arriver les \omterm{def:mx:the-limit-order-book:orders}{ordres à cours limité}, les \omterm{def:mx:the-limit-order-book:orders}{ordres au marché} et les annulations
sous forme de processus de Poisson indépendants dont les intensités peuvent dépendre de la distance aux meilleurs prix, sans qu'aucun intervenant ne choisisse quoi que ce soit~; sa fourchette,
sa profondeur, sa volatilité et ses probabilités d'exécution découlent des seules intensités.
\end{definition}

Smith, Farmer, Gillemot et Krishnamurthy (2003) ont construit le modèle avec des \omterm{def:mx:the-limit-order-book:orders}{ordres à cours limité} placés au hasard sur une plage de prix, des \omterm{def:mx:the-limit-order-book:orders}{ordres au marché} et
des annulations, et ont établi par analyse dimensionnelle comment la fourchette, la profondeur, la volatilité et l'impact varient avec les intensités~; toutes les intensités étant mesurables
dans les données, le modèle n'a aucun paramètre libre. Cont, Stoikov et Talreja (2010) ont écrit le carnet comme un ensemble de files de naissance et de mort (livre~4, chapitre~8) avec
des annulations à un taux proportionnel à la taille de la file, l'ont ajusté aux données de la Tokyo Stock Exchange et calculé, parmi d'autres probabilités conditionnelles,
celle d'exécuter un ordre au cours acheteur avant que le cours vendeur ne bouge. Le résultat suivant en constitue la brique de base.

\begin{proposition}[Probabilité d'exécution derrière $n$ ordres]\label{prop:mx:the-limit-order-book:fill}
Un ordre rejoint un \omterm{def:mx:the-limit-order-book:tick}{niveau de prix} derrière $n$ ordres d'une unité chacun. Les \omterm{def:mx:the-limit-order-book:orders}{ordres au marché} arrivent au taux $\mu$ et chacun prélève une unité en tête de file~; chaque
ordre qui le précède est annulé au taux $\theta$~; l'ordre lui-même est annulé au taux $\nu$. Toutes les horloges sont indépendantes et exponentielles. La probabilité
qu'il soit exécuté avant d'être annulé est
\[
P_n=\frac{\mu}{\mu+\nu}\prod_{k=1}^{n}\frac{\mu+k\theta}{\mu+k\theta+\nu},
\]
et, lorsque $\nu=0$, l'espérance du temps jusqu'à l'exécution vaut $1/\mu+\sum_{k=1}^{n}1/(\mu+k\theta)$.
\end{proposition}

\begin{proof}
Avec $k\ge1$ ordres devant, le prochain événement est un départ devant (taux $\mu+k\theta$) ou notre annulation (taux $\nu$)~; par absence de mémoire, le
départ survient en premier avec la probabilité $(\mu+k\theta)/(\mu+k\theta+\nu)$, et la file repart alors de $k-1$. En tête ($k=0$), le prochain
\omterm{def:mx:the-limit-order-book:orders}{ordre au marché} nous exécute avec la probabilité $\mu/(\mu+\nu)$. Les étapes sont indépendantes, donc les probabilités se multiplient. Avec $\nu=0$, les
durées sont exponentielles de moyennes $1/(\mu+k\theta)$ et $1/\mu$~; l'espérance totale est leur somme.
\end{proof}

\omcode{code/microstructure/01-the-limit-order-book/python/mx_lob.py}{115}{121}{La proposition en code.}\label{lst:mx:the-limit-order-book:fill}

La proposition vaut dans une file où personne ne passe devant. Dans un carnet réel, comme dans le \omterm{def:mx:the-limit-order-book:zi}{modèle à intelligence nulle}, quelqu'un le fait~: un acheteur qui place
un cours acheteur d'un tick au-dessus du nôtre devient le meilleur cours acheteur, et les ventes au marché lui sont adressées. La \Cref{fig:mx:the-limit-order-book:fill} mesure l'écart. Le
marché à intelligence nulle du chapitre (ordres unitaires, \omterm{def:mx:the-limit-order-book:orders}{ordres à cours limité} au taux de 0.5 par prix et par côté sur 20 prix, \omterm{def:mx:the-limit-order-book:orders}{ordres au marché} au taux de 2 par côté, annulations au taux de
0.05 par ordre) suit chaque ordre au cours acheteur qui rejoint l'arrière du meilleur cours acheteur. Avec l'amélioration de prix désactivée, ses taux d'exécution concordent avec ceux de la proposition, avec
$\mu=2$, $\theta=\nu=0.05$~: derrière cinq ordres, 0.871 mesuré contre 0.870. Avec l'amélioration de prix activée, comme dans le modèle d'origine, ils tombent à 0.672 derrière
cinq ordres et à 0.535 derrière huit~: entre un quart et un tiers des exécutions que promet la proposition sont prises par des ordres qui arrivent
plus tard à un meilleur prix. Désactiver
l'amélioration de prix a un coût propre~: la fourchette, de 2.55 ticks avec elle, croît sans limite sans elle, car rien ne la resserre jamais.
C'est l'amélioration de prix qui maintient la fourchette finie et qui fait que la position dans la file n'est pas une garantie.

\begin{omfigure}
\begin{tikzpicture}
\begin{axis}[name=fill,width=7.2cm,height=5.4cm,title={\small probabilité d'exécution},xlabel={\small ordres devant $n$},
  tick label style={font=\footnotesize},xtick={1,2,3,4,5,6,7,8},ymin=0.4,ymax=1.0,grid=major,grid style={gray!25},
  legend columns=1,legend style={font=\footnotesize,at={(0.5,-0.3)},anchor=north,draw=none,fill=none},table/col sep=comma]
\addplot[black,thick] table[x=n,y=formula]{figdata/microstructure/01-the-limit-order-book/fill.csv};
\addplot[omDef,thick,mark=o,only marks] table[x=n,y=zi_noimprove]{figdata/microstructure/01-the-limit-order-book/fill.csv};
\addplot[omThm,thick,mark=square*] table[x=n,y=zi]{figdata/microstructure/01-the-limit-order-book/fill.csv};
\legend{proposition ($\mu=2$; $\theta=\nu=0.05$),modèle à intelligence nulle sans amélioration,modèle à intelligence nulle avec amélioration}
\end{axis}
\begin{axis}[at={(fill.east)},anchor=west,xshift=1.3cm,width=7.2cm,height=5.4cm,title={\small profondeur affichée moyenne (actions)},
  xlabel={\small ticks depuis le meilleur},tick label style={font=\footnotesize},xtick={0,1,2,3,4,5,6,7,8,9},ymin=0,ymax=1700,grid=major,
  grid style={gray!25},ybar,bar width=4pt,legend columns=2,legend style={font=\footnotesize,at={(0.5,-0.3)},anchor=north,draw=none,fill=none},
  table/col sep=comma]
\addplot[fill=omDef!60,draw=omDef] table[x=k,y=bid]{figdata/microstructure/01-the-limit-order-book/depth.csv};
\addplot[fill=omThm!60,draw=omThm] table[x=k,y=ask]{figdata/microstructure/01-the-limit-order-book/depth.csv};
\legend{cours acheteur,cours vendeur}
\end{axis}
\end{tikzpicture}

\omcaption{À gauche~: la probabilité qu'un ordre placé à l'arrière du meilleur cours acheteur soit exécuté avant d'être annulé, en fonction des ordres qui le précèdent~: la
proposition, et le marché à intelligence nulle avec et sans amélioration de prix (400\,000 événements chacun~; au-delà de huit ordres devant
les ordres suivis sont trop peu nombreux, moins d'une centaine par point). À droite~: la profondeur affichée moyenne à chaque
distance du meilleur prix dans l'heure simulée de \texttt{firm.tape}, échantillonnée toutes les secondes~; elle culmine à un tick derrière le meilleur. Données~:
\texttt{mx\_lob.zi\_market}, \texttt{mx\_lob.depth\_profile}.}\label{fig:mx:the-limit-order-book:fill}
\end{omfigure}

Le panneau de droite de la \cref{fig:mx:the-limit-order-book:fill} montre une forme que le chapitre~3 retrouvera dans les ouvrages réels~: la profondeur affichée n'est pas maximale au meilleur
cours, mais à un tick derrière lui, car c'est au meilleur niveau que se concentrent les exécutions et la plupart des annulations.

\section{Tutoriel~: reconstruire un carnet et chronométrer sa file}\label{tut:mx:the-limit-order-book:tut}

\begin{tutorial}
\textbf{Objectif.} Reconstruire le carnet de l'heure simulée dans deux structures, mesurer sa répartition des événements et ses files, et tester la
proposition sur la probabilité d'exécution dans un marché à intelligence nulle. \textbf{État final~:} les chiffres de la section~2, les deux structures concordant, et
\cref{fig:mx:the-limit-order-book:fill}.
\begin{tutsteps}
\item \textbf{Reconstruire.} \texttt{mx\_lob.tree\_vs\_array(session())} rejoue 31\,535 messages dans un \texttt{MessageBook} et un \texttt{ArrayBook} et
compare leurs dix meilleurs niveaux tous les dix messages~: 3\,154 comparaisons, aucune divergence.
\item \textbf{Mesurer.} \texttt{event\_mix}, \texttt{best\_survival} et \texttt{joiners} donnent les parts des ajouts, des annulations et des exécutions, les durées de vie
des ordres, la durée pendant laquelle un meilleur cours tient et ce que deviennent les ordres qui rejoignent la meilleure file.
\item \textbf{File.} Évaluer \texttt{fill\_probability} et la comparer à \texttt{queue\_mc}~; lancer \texttt{zi\_market(improve=False)} et
\texttt{zi\_market(improve=True)}~; tracer la figure avec \texttt{fig\_lob.py}.
\item \textbf{Jumeaux.} Lancer \texttt{make\_lob\_fixture.py}~; construire \texttt{cpp/firm\_lob\_test.cpp} et \texttt{cargo test} dans \texttt{rust/}~: les deux reproduisent
les 147 instantanés du carnet Python.
\end{tutsteps}
\textbf{Pour aller plus loin.} Donner aux ordres du marché à intelligence nulle des tailles tirées d'une loi géométrique et observer de combien la proposition, comptée en
ordres, s'écarte du taux d'exécution~; diviser par deux le taux d'annulation et observer l'évolution de la fourchette et de la profondeur au meilleur cours.
\end{tutorial}

\section{Construction~: le carnet d'ordres de référence}\label{bld:mx:the-limit-order-book:lob}

\begin{build}
\bfield{Rôle} Le carnet dans lequel chaque composant ultérieur de la firme stocke ses ordres~: le moteur du simulateur de bourse (chapitre~26), les vues propres
des participants sur le marché, et les tests des carnets de production du livre~13.
\bfield{Interface} \texttt{Order(ref, side, price, qty, visible, seq, owner)}~; \texttt{LimitOrderBook}~: \texttt{add}, \texttt{remove}, \texttt{reduce},
\texttt{get}, \texttt{best}, \texttt{best\_visible}, \texttt{prices}, \texttt{level\_orders}, \texttt{top}, \texttt{depth}, \texttt{check}~;
\texttt{MessageBook.apply(kind, ref, side, price, qty, new\_ref)} pour A, X, D, E, C et U~; \texttt{ArrayBook(lo, hi, tick)}~; C++20
\texttt{firm::lob::MessageBook} et Rust \texttt{firm\_lob::MessageBook} avec le même \texttt{l2\_lines}.
\bfield{Règles} Les ordres affichés passent avant les ordres cachés, puis le temps prime à chaque prix~; une réduction conserve la place~; un remplacement arrive sous une nouvelle référence à la fin de la file~;
le carnet enregistre et ne fait jamais d'appariement~; \texttt{check()} vérifie que chaque ordre se trouve exactement dans son niveau.
\bfield{Tests d'acceptation} \texttt{code/firm/lob/tests/}~: la priorité vérifiée à la main~; les deux structures Python concordant sur une session simulée~; le remplacement~;
les 147 instantanés de niveau 2 du fixture en C++20 et en Rust.
\bfield{Pour aller plus loin} Une liste chaînée intrusive par niveau pour une suppression en $O(1)$~; un carnet en tableau qui se recentre lorsque le prix sort de sa bande.
\end{build}

\begin{omsources}
\item M.~D.~Gould, M.~A.~Porter, S.~Williams, M.~McDonald, D.~J.~Fenn et S.~D.~Howison, ``Limit order books'', \emph{Quantitative Finance}
13(11), 2013.
\item E.~Smith, J.~D.~Farmer, L.~Gillemot et S.~Krishnamurthy, ``Statistical theory of the continuous double auction'', \emph{Quantitative
Finance} 3(6), 2003.
\item R.~Cont, S.~Stoikov et R.~Talreja, ``A stochastic model for order book dynamics'', \emph{Operations Research} 58(3), 2010.
\item Nasdaq, spécification \emph{TotalView-ITCH 5.0}~; CME Group, documentation des systèmes clients, ``Order functionalities''.
\end{omsources}

\section{Exercices}

\begin{exercise}[$\star$]\label{exo:mx:the-limit-order-book:1}
Le meilleur cours acheteur est de 100.00 avec 700 actions réparties en trois ordres de 300, 200 et 200~; le meilleur cours vendeur est de 100.01. Un ordre de vente au marché de 450 actions arrive. Quelles transactions
se produisent, et quel est le nouveau niveau 1~?
\end{exercise}

\begin{exercise}[$\star$]\label{exo:mx:the-limit-order-book:2}
Un ordre rejoint derrière trois ordres unitaires~; les \omterm{def:mx:the-limit-order-book:orders}{ordres au marché} arrivent au rythme de 1 par seconde, chaque ordre qui le précède s'annule à 0.1 par seconde, et le propriétaire de l'ordre
l'annule à 0.2 par seconde. Quelle est sa probabilité d'exécution~?
\end{exercise}

\begin{exercise}[$\star$]\label{exo:mx:the-limit-order-book:3}
Un trader réduit d'abord un ordre de 500 à 300 actions, puis le remonte à 500. Qu'advient-il de sa position dans la file à chaque étape sur CME Globex,
et comment chaque étape apparaît-elle dans un flux ITCH~?
\end{exercise}

\begin{exercise}[$\star\star$]\label{exo:mx:the-limit-order-book:4}
Un tableau indexé par le prix couvre les prix situés à moins de 1~\% de 100.00 pour une action dont le tick vaut un centime. Combien d'emplacements faut-il, et que se passe-t-il
un jour où l'action varie de 3~\%~?
\end{exercise}

\begin{exercise}[$\star\star$]\label{exo:mx:the-limit-order-book:5}
Sans annulation de sa part, un ordre attend derrière dix ordres unitaires~; les \omterm{def:mx:the-limit-order-book:orders}{ordres au marché} arrivent à 0.5 par seconde et chaque ordre qui le précède s'annule à 0.02 par
seconde. Quelle est son attente moyenne~?
\end{exercise}

\begin{exercise}[$\star\star$]\label{exo:mx:the-limit-order-book:6}
Pourquoi la profondeur affichée dans la \cref{fig:mx:the-limit-order-book:fill} est-elle plus grande à un tick derrière le meilleur cours qu'au meilleur cours~?
\end{exercise}

\begin{exercise}[$\star\star\star$]\label{exo:mx:the-limit-order-book:7}
\emph{Programmation.} Relancer \texttt{zi\_market} avec des \omterm{def:mx:the-limit-order-book:orders}{ordres au marché} à un rythme de 1 au lieu de 2. Qu'advient-il de la fourchette, et de l'écart entre la
proposition et la probabilité d'exécution mesurée derrière cinq ordres~? Expliquer les deux points.
\end{exercise}

\begin{exercise}[$\star\star\star$]\label{exo:mx:the-limit-order-book:8}
\emph{Trouver la faille.} «~Notre ordre est en tête du meilleur cours acheteur depuis dix secondes, donc d'après la proposition il sera exécuté avec la probabilité
$\mu/(\mu+\nu)$ quoi qu'il arrive ensuite.~»
\end{exercise}

\section{Problème~: la file au meilleur cours acheteur}

\begin{problem}[{Problème du week-end --- la file au meilleur cours acheteur}]\label{pb:mx:the-limit-order-book:1}
Un desk demande combien de temps ses ordres passifs attendent et combien d'entre eux sont exécutés. On dispose de l'heure simulée de \texttt{firm.tape} et du marché
à intelligence nulle du chapitre.

\medskip\noindent\textbf{Partie I --- le flux.}
\begin{enumerate}
\item Combien de messages contient l'heure, et dans quelles proportions d'ajouts, d'annulations et d'exécutions~?
\item Combien de messages correspondent à chaque transaction, et quelle part des ordres sort du carnet par une exécution~?
\item Combien de temps vit l'ordre médian, et un ordre exécuté vit-il plus longtemps qu'un ordre annulé~?
\item Combien de temps un meilleur cours acheteur tient-il, dans le cas médian~?
\end{enumerate}

\medskip\noindent\textbf{Partie II --- la file au meilleur cours.}
\begin{enumerate}[resume]
\item Combien d'ordres rejoignent l'arrière du meilleur cours acheteur ou vendeur, et combien d'actions trouvent-ils devant eux~?
\item Quelle part d'entre eux finit par être exécutée, et quelle part l'est en totalité~?
\item Où la profondeur affichée moyenne culmine-t-elle, et pourquoi à cet endroit~?
\item Lesquels de ces chiffres changeraient si les annulations étaient plus lentes~?
\end{enumerate}

\medskip\noindent\textbf{Partie III --- le modèle.}
\begin{enumerate}[resume]
\item Énoncer la proposition sur la probabilité d'exécution et ses hypothèses.
\item Calculer $P_1$, $P_5$ et $P_{12}$ pour $\mu=2$, $\theta=\nu=0.05$.
\item Vérifier $P_5$ par simulation de la file stylisée.
\item Quel est le temps d'attente espéré derrière cinq ordres lorsque $\nu=0$~?
\end{enumerate}

\medskip\noindent\textbf{Partie IV --- le marché.}
\begin{enumerate}[resume]
\item Quels taux d'exécution le \omterm{def:mx:the-limit-order-book:zi}{modèle à intelligence nulle} donne-t-il sans amélioration de prix, et pourquoi concordent-ils avec la proposition~?
\item Que donne-t-il avec amélioration, derrière un, cinq et huit ordres~?
\item Énoncer le \emph{résultat nommé}~: la probabilité d'exécution d'un ordre placé à l'arrière du meilleur cours acheteur face aux ordres qui le précèdent, et la part de celle-ci
que l'amélioration de prix lui enlève.
\item Pourquoi le marché sans amélioration présente-t-il une fourchette non bornée~?
\item Quelles sont la fourchette et la profondeur moyenne au meilleur niveau du \omterm{def:mx:the-limit-order-book:zi}{modèle à intelligence nulle} avec amélioration~?
\item Quelle hypothèse de la proposition est la plus souvent mise en défaut dans un carnet d'ordres réel, et dans quel sens biaise-t-elle la réponse~?
\item Que faudrait-il mesurer dans des données réelles pour estimer $\mu$, $\theta$ et $\nu$~?
\item En une phrase~: que vaut une position dans la file, avant que le chapitre~6 ne la valorise~?
\end{enumerate}
\end{problem}

\section{Questions d'entretien}

\begin{interviewq}[$\star$ \iqroles{développeur}]\label{iq:mx:the-limit-order-book:1}
Concevez les structures de données d'un \omterm{def:mx:the-limit-order-book:lob}{carnet d'ordres à cours limité}. Quels sont les coûts d'ajout, d'annulation, d'exécution et de recherche du meilleur cours~?
\end{interviewq}

\begin{interviewq}[$\star$ \iqroles{trader, chercheur}]\label{iq:mx:the-limit-order-book:2}
Pourquoi augmenter la taille d'un ordre lui fait-il perdre sa position dans la file, alors que la diminuer ne le fait pas~?
\end{interviewq}

\begin{interviewq}[$\star\star$ \iqroles{chercheur}]\label{iq:mx:the-limit-order-book:3}
Vous êtes cinquième dans la file. Des \omterm{def:mx:the-limit-order-book:orders}{ordres au marché} arrivent au taux $\mu$, et chacun des ordres qui vous précèdent s'annule au taux $\theta$. Quelle est la probabilité que vous soyez exécuté avant de
vous annuler au taux $\nu$~?
\end{interviewq}

\begin{interviewq}[$\star\star$ \iqroles{développeur}]\label{iq:mx:the-limit-order-book:4}
Un carnet à tableau indexé par les prix est rapide. Quand échoue-t-il, et que faites-vous pour y remédier~?
\end{interviewq}

\begin{interviewq}[$\star\star$ \iqroles{chercheur, trader}]\label{iq:mx:the-limit-order-book:5}
La plupart des ordres d'un carnet moderne sont annulés. Est-ce un problème pour le marché~?
\end{interviewq}

\begin{interviewq}[$\star\star\star$ \iqroles{chercheur}]\label{iq:mx:the-limit-order-book:6}
Que reproduit correctement un \omterm{def:mx:the-limit-order-book:zi}{modèle à intelligence nulle} du carnet d'ordres, et que rate-t-il~?
\end{interviewq}
